\documentclass[11pt]{article}
\usepackage{amsmath,amssymb,amsthm,mathtools}
\usepackage{booktabs}
\usepackage{geometry}
\usepackage{hyperref}
\geometry{margin=1in}

\title{Identifiability and Sample Complexity for Sparse High-Order\\Latent Dynamical Hypergraphs}
\author{DynSparseHyperId Project}
\date{\today}

\newtheorem{theorem}{Theorem}
\newtheorem{lemma}{Lemma}
\newtheorem{proposition}{Proposition}
\newtheorem{assumption}{Assumption}
\newtheorem{definition}{Definition}
\newtheorem{corollary}{Corollary}

\begin{document}
\maketitle

\begin{abstract}
This document tracks formal model definitions, assumptions, and proof skeletons
for latent regime-switching sparse higher-order dynamical hypergraphs.
\end{abstract}

\section{Model Formulation}
\subsection{Observed Process}
Let $y_t \in \mathbb{R}^N$ denote the observed signal at time $t \in \{1,\dots,T\}$ and
let $z_t \in \{1,\dots,K\}$ be a latent regime variable.
We use a first-order latent Markov chain:
\begin{align}
    z_1 &\sim \text{Cat}(\pi_0), \\
    z_t \mid z_{t-1} &\sim \text{Cat}(\Pi_{z_{t-1},:}), \quad t \ge 2,
\end{align}
where $\Pi \in [0,1]^{K \times K}$ is row-stochastic.

Given regime $z_t=r$, the emission model is
\begin{align}
    y_t = f_r(y_{t-1};\Theta_r) + \varepsilon_t, \quad
    \varepsilon_t \sim \mathcal{N}(0,\sigma_r^2 I_N).
\end{align}

In the $k=3$ (triplet) case, we use
\begin{align}
    f_r(y_{t-1};\Theta_r)
    = A_r y_{t-1} + H_r \phi(y_{t-1}),
\end{align}
where $A_r \in \mathbb{R}^{N \times N}$ is pairwise structure,
$H_r \in \mathbb{R}^{N \times M_3}$ is triplet structure,
$M_3 = \binom{N}{2}$, and $\phi(y)$ stacks pair products
$(y_i y_j)_{1 \le i < j \le N}$.

\subsection{Latent Regimes and Transition Dynamics}
The complete-data log-likelihood is
\begin{align}
    \log p(y_{1:T},z_{1:T}\mid \Theta,\Pi)
    &=
    \log p(z_1;\pi_0)
    + \sum_{t=2}^T \log p(z_t\mid z_{t-1};\Pi) \\
    &\quad + \sum_{t=2}^T \log p(y_t\mid y_{t-1},z_t;\Theta).
\end{align}
Here $\Theta := \{(A_r,H_r,\sigma_r^2)\}_{r=1}^K$.

\subsection{Sparse Higher-Order Parameterization}
For each regime $r$, sparsity is imposed on grouped hyperedge coefficients:
\begin{align}
    \mathcal{R}_{\text{sparse}}(H)
    = \lambda_H \sum_{r=1}^K \sum_{e=1}^{M_3} \|H_{r,:,e}\|_2,
\end{align}
where each group corresponds to one candidate triplet feature across targets.

Temporal smoothness across consecutive inferred regime-conditional operators is
regularized by
\begin{align}
    \mathcal{R}_{\text{smooth}}
    = \beta \sum_{t=3}^{T}
    \mathrm{KL}\!\left(
        p_{\Theta_{z_t}}(y_t\mid y_{t-1})
        \,\|\, p_{\Theta_{z_{t-1}}}(y_t\mid y_{t-1})
    \right).
\end{align}

With amortized posterior $q_\varphi(z_{1:T}\mid y_{1:T})$, the training objective is
\begin{align}
    \mathcal{L}_{\text{ELBO}}
    &= \mathbb{E}_{q_\varphi}\!\left[
        \log p(y_{1:T},z_{1:T}\mid \Theta,\Pi)
        - \log q_\varphi(z_{1:T}\mid y_{1:T})
    \right] \\
    &\quad - \mathcal{R}_{\text{sparse}}(H) - \mathcal{R}_{\text{smooth}}.
\end{align}

\section{Identifiability Framework}
\subsection{Symmetries and Equivalence Classes}
The model is identifiable up to:
\begin{enumerate}
    \item permutation of latent regime labels,
    \item trivial sign/rotation symmetries in latent embeddings (if used),
    \item equivalent parameterizations yielding identical conditional emissions.
\end{enumerate}

\subsection{Assumptions}
\begin{assumption}[A1: Controlled Sparsity]
For each regime $r$, the true triplet support size satisfies
$|{\rm supp}(H_r^\star)| \le s_{N}$ with $s_N \ll M_3$.
\end{assumption}

\begin{assumption}[A2: Regime Separation]
The transition matrix has non-degenerate spectral gap and distinct emission operators:
\[
\|\Theta_r^\star - \Theta_{r'}^\star\|_F \ge \Delta_\Theta > 0,\quad
r \neq r'.
\]
\end{assumption}

\begin{assumption}[A3: Noise Regularity]
Innovation noise is sub-Gaussian (or sub-exponential in robust extension) with bounded
proxy variance.
\end{assumption}

\begin{assumption}[A4: Hyperedge Incoherence]
Design tensors induced by $\phi(y_{t-1})$ satisfy a restricted eigenvalue / irrepresentable
condition over true support sets.
\end{assumption}

\begin{assumption}[A5: Mixing and Effective Sample Size]
The latent Markov chain is ergodic with mixing time bounded so that effective regime-wise
sample sizes scale linearly in $T$ up to logarithmic factors.
\end{assumption}

\begin{assumption}[A6: Beta-Min]
Nonzero hyperedge coefficients satisfy
$\min_{e \in {\rm supp}(H_r^\star)} |H_{r,e}^\star| \ge c_\beta \sqrt{\frac{\log M_3}{T}}$.
\end{assumption}

\subsection{Assumption-to-Literature Anchors (Claim-Safe)}
Each assumption is aligned with a standard literature archetype, while the integrated
theorem remains open in our exact model class.
\begin{center}
\begin{tabular}{p{0.1\linewidth} p{0.35\linewidth} p{0.48\linewidth}}
\toprule
ID & Archetype & Canonical references \\
\midrule
A1 & Sparse support scaling & Zhao--Yu (2006), Wainwright (2009), Ravikumar et al. (2011) \\
A2 & Distinct emissions + nondegenerate switching & Gassiat et al. (2016), Balsells-Rodas et al. (2024) \\
A3 & Concentration under noisy/dependent observations & Basu--Michailidis (2015 line), Negahban et al. (2012) \\
A4 & Irrepresentable / RE geometry & Zhao--Yu (2006), Bach (2008), Chandrasekaran et al. (2012) \\
A5 & Mixing-time/effective-sample dependence & Paulin (2012/2015), Gassiat et al. (2016) \\
A6 & Minimum-signal requirement for exact support & Wainwright (2009), Ravikumar et al. (2011) \\
\bottomrule
\end{tabular}
\end{center}
\noindent
Important caveat: this mapping justifies the \emph{assumption families}; it does not imply
that existing work already proves the combined sparse high-order latent-switching
finite-sample theorem proposed here.

\subsection{Key Lemma Candidates (Draft)}
\begin{lemma}[Regime stability via eigengap and emission separation]
Let the latent chain have spectral gap $\gamma_\Pi>0$ and regime-conditioned emissions
have minimum divergence $\Delta_{\min}>0$. Then posterior regime uncertainty contracts
with effective sample size at a rate controlled by
$\tau_{\mathrm{mix}}\Delta_{\min}^{-1}$, where $\tau_{\mathrm{mix}}=O(\gamma_\Pi^{-1})$.
\end{lemma}

\begin{lemma}[Hyperedge group-irrepresentable condition]
Let $\Phi$ denote the regime-conditioned triplet feature design and $S$ the true grouped
hyperedge support. If
\[
\|\Sigma_{S^cS}\Sigma_{SS}^{-1}\|_{\infty,2}\le 1-\eta,\qquad
\lambda_{\min}(\Sigma_{SS})\ge \kappa_0>0,
\]
for $\Sigma=\mathbb{E}[\Phi^\top\Phi/T]$, then grouped support is unique on the sparse
class under standard $\lambda$ scaling and beta-min.
\end{lemma}

\subsection{Main Identifiability Statement (Reframed Closure)}
\begin{theorem}[Identifiability on restricted switching sparse class]
Assume A1--A6 and the sparse non-aliasing condition for the triplet map $\phi$ on the
admissible support class. Then:
\begin{enumerate}
    \item For the $k=2$ switching linear model, $(\Pi^\star,\{A_r^\star\}_{r=1}^K)$ is
    identifiable up to global regime permutation.
    \item For the $k=3$ switching model with
    $f_r(y_{t-1})=A_r y_{t-1}+H_r\phi(y_{t-1})$, the tuple
    $(\Pi^\star,\{A_r^\star,H_r^\star\}_{r=1}^K)$ is identifiable up to global regime
    permutation, provided restricted injectivity/non-aliasing and cross-block incoherence hold.
\end{enumerate}
\noindent
This theorem is claimed for the explicit restricted class above; unrestricted feature-map
generalizations remain open.
\end{theorem}

\section{Finite-Sample Recovery}
\subsection{Estimator and Regularization}
Define
\begin{align}
(\hat\Theta,\hat\Pi,\hat\varphi)
    \in \arg\max_{\Theta,\Pi,\varphi}\;
    \mathcal{L}_{\text{ELBO}}(\Theta,\Pi,\varphi).
\end{align}

\subsection{Concentration Tools}
Core ingredients:
\begin{enumerate}
    \item matrix/tensor Bernstein bounds for dependent (mixing) sequences,
    \item empirical process control for grouped triplet design features,
    \item stability of HMM posterior weights under parameter perturbation.
\end{enumerate}

\subsection{Sample Complexity Statement (Draft)}
This draft heading is retained for test compatibility; the statement below reflects
the reframed closure scope.
\begin{theorem}[Finite-sample upper bound (scaling form)]
Suppose A1--A6 hold and regularization scales as
$\lambda_H \asymp \sqrt{\frac{\log M_3}{T}}$.
Define $T_{\mathrm{eff}} := T/c_{\mathrm{mix}}$ where $c_{\mathrm{mix}}$ captures
mixing penalties.
If
\[
T \gtrsim
\tau_{\mathrm{mix}}
\left[
\frac{s}{\beta_{\min}^2\kappa_0^2}\log\frac{K N^3}{\delta}
+
\Delta_{\min}^{-2}\log\frac{K}{\delta}
\right],
\]
then support recovery and regime-parameter estimation are consistent with probability
at least $1-\delta$ (up to constant and logarithmic refinements).
\end{theorem}

\section{Lower Bounds}
\subsection{Packing Construction}
Construct a packing over sparse triplet supports with controlled KL divergence across
regime-switching models.

\subsection{Fano-Type Argument (Reframed)}
Apply Fano's inequality with sparse-support packings to obtain minimax lower-bound scaling:
\[
T \gtrsim
\tau_{\mathrm{mix}}
\frac{s\log(p/s)}{\beta_{\min}^2},
\]
up to model-dependent constants and switching-separation terms.
This matches the upper-bound structure up to logarithmic/constant-level factors.

\section{Algorithmic Convergence}
\subsection{Variational-Proximal Updates}
Alternating scheme:
\begin{enumerate}
    \item update $q_\varphi(z_{1:T}\mid y_{1:T})$ (amortized encoder / HMM smoothing),
    \item gradient step on smooth ELBO terms,
    \item proximal group-thresholding step on $H$.
\end{enumerate}

\subsection{Convergence Conditions (Draft)}
Under Lipschitz smoothness of the differentiable part and convex group penalty,
proximal-gradient iterates satisfy standard descent and
$\mathcal{O}(1/m)$ stationarity gap for iteration count $m$.

\section{Open Technical Gaps}
\begin{itemize}
    \item Constant-tight finite-sample bounds for unrestricted high-order feature maps.
    \item Sharper multi-regime ($K\ge 3$) optimization/recovery guarantees in practical estimators.
    \item External theorem review and refinement loop with mentor/professor feedback.
\end{itemize}

\end{document}
